\chapter{实验设计与结果组织}

\section{实验章节定位}
本章作为全文倒数第二章，统一承接前文的 benchmark 构建、方法设计与结果验证任务。与仅对少量场景做“点状展示”的写法不同，本文当前版本明确以
\path{benchmark_88_2clusters_11scenes_8methods_fixedfold.yaml}
对应的 88-run fixed-fold benchmark 作为主实验口径。也就是说，本章首先围绕当前必须完成的正式 benchmark 给出实验组织、结果表和图件占位；至于更细粒度的模块级消融，则在本章后部作为后续补充实验计划保留结构化位置。

这种组织方式的好处在于：一方面，论文主结果与当前实际要跑的实验配置完全一致，避免正文口径与代码配置脱节；另一方面，第三章至第六章已经把 RCTA 的创新模块分章展开，即使暂时还没有逐模块消融结果，论文的技术逻辑仍然完整，后续只需在既有位置回填补充实验。

\section{数据集与 benchmark 场景}
\subsection{TEP 领域自适应数据集}
本文采用 Tennessee Eastman Process Domain Adaptation 数据集开展实验。该数据集将 6 个操作模式视为 6 个不同领域，标签空间为 29 类，即正常工况和 28 类故障。与普通单域分类任务相比，该数据集更适合用于分析工业过程在工况切换条件下产生的分布偏移问题，并观察不同领域自适应机制对多类别工业时序诊断的影响~\cite{montesuma2023benchmark}。

根据当前工作区中的 \texttt{data/benchmark/manifest.json}，6 个域的样本规模与默认 Fold 1 划分情况如表~\ref{tab:dataset-stats} 所示。可以看到，各域样本总量均接近 2900，默认采用 Fold 1 作为评估折时，训练样本量约为 2300，评估样本量约为 570。这样的规模既能够支撑本文的系统化 benchmark 比较，也保留了工业时序诊断任务的复杂性。

\begin{table}[htbp]
\centering
\caption{TEP 领域自适应数据集的域规模与默认 Fold 1 划分}
\label{tab:dataset-stats}
\begin{tabular}{lccccc}
\hline
域 & 样本总数 & 训练样本数 & 评估样本数 & 单样本形状 & 类别数 \\
\hline
mode1 & 2900 & 2320 & 580 & 600$\times$34 & 29 \\
mode2 & 2845 & 2278 & 567 & 600$\times$34 & 29 \\
mode3 & 2899 & 2320 & 579 & 600$\times$34 & 29 \\
mode4 & 2865 & 2294 & 571 & 600$\times$34 & 29 \\
mode5 & 2883 & 2307 & 576 & 600$\times$34 & 29 \\
mode6 & 2897 & 2318 & 579 & 600$\times$34 & 29 \\
\hline
\end{tabular}
\end{table}

\subsection{11 个场景的 fixed-fold benchmark}
当前 benchmark 不追求覆盖全部 30 个有向单源场景，而是采用“两簇 + 代表性多源场景”的设计。具体而言，一簇为 mode3 与 mode6 的互迁，另一簇为 mode1、mode2、mode5 之间的互迁，并在后者基础上补充 3 个多源到单目标场景。这样得到总计 11 个迁移场景，其中 8 个为单源场景，3 个为多源场景。

表~\ref{tab:benchmark-scenes} 给出了本 benchmark 的场景构成与固定折分配。固定折策略的作用在于：所有方法在同一场景上共享相同的源折与目标折，从而避免方法之间由于折采样差异而产生额外噪声，使比较尽量聚焦于方法本身。

\begin{table}[htbp]
\centering
\caption{88-run fixed-fold benchmark 的 11 个场景与折分配}
\label{tab:benchmark-scenes}
\small
\resizebox{\textwidth}{!}{%
\begin{tabular}{lllll}
\hline
编号 & 场景类型 & 迁移场景 & 源折设置 & 目标折 \\
\hline
S1 & 单源 & mode3$\rightarrow$mode6 & Fold 1 & Fold 1 \\
S2 & 单源 & mode6$\rightarrow$mode3 & Fold 1 & Fold 5 \\
S3 & 单源 & mode1$\rightarrow$mode2 & Fold 2 & Fold 4 \\
S4 & 单源 & mode2$\rightarrow$mode1 & Fold 5 & Fold 3 \\
S5 & 单源 & mode1$\rightarrow$mode5 & Fold 4 & Fold 3 \\
S6 & 单源 & mode5$\rightarrow$mode1 & Fold 4 & Fold 4 \\
S7 & 单源 & mode2$\rightarrow$mode5 & Fold 1 & Fold 3 \\
S8 & 单源 & mode5$\rightarrow$mode2 & Fold 3 & Fold 4 \\
M1 & 多源 & mode1+mode2$\rightarrow$mode5 & mode1: Fold 4, mode2: Fold 1 & Fold 3 \\
M2 & 多源 & mode1+mode5$\rightarrow$mode2 & mode1: Fold 2, mode5: Fold 3 & Fold 4 \\
M3 & 多源 & mode2+mode5$\rightarrow$mode1 & mode2: Fold 4, mode5: Fold 4 & Fold 4 \\
\hline
\end{tabular}
}
\normalsize
\end{table}

在实验数量上，当前 benchmark 共包含 8 种方法和 11 个场景，因此总运行次数为
\begin{equation}
8\times 11=88.
\end{equation}
这一规模既明显大于只比较个别迁移方向的“小样本结果”，又仍然处于本科论文阶段可以组织与管理的范围内。

\section{统一实验协议}
\subsection{输入预处理与归一化}
数据集中的每个样本均表示为固定长度的多变量时间窗口，可形式化写为
\begin{equation}
x\in\mathbb{R}^{T\times C},
\end{equation}
其中 $T=600$，$C=34$。在进入神经网络前，数据被转置为 channels-first 形式：
\begin{equation}
x^\prime=\mathrm{Transpose}(x)\in\mathbb{R}^{C\times T}.
\end{equation}
为保证不同方法具有可比性，本文采用域内标准化，即对每个域分别基于该域全量样本统计量执行标准化：
\begin{equation}
\tilde{x}_{n,t,c}=\frac{x_{n,t,c}-\mu_c}{\sigma_c+\varepsilon}.
\end{equation}
其中 $\mu_c$ 和 $\sigma_c$ 为该域第 $c$ 个通道的均值和标准差。该策略对应当前数据配置中的 \path{normalization_scope=domain}。

\subsection{方法集合与训练设置}
当前 benchmark 涉及的 8 种方法为 Source Only、Target Only、CORAL、DAN、DANN、CDAN、DeepJDOT 和 RCTA。
其中，Target Only 使用目标域标签进行监督训练，因此在严格意义上不属于无监督领域自适应方法，而是作为监督参考上界参与比较。其余 7 种方法中，Source Only 作为无自适应对照组，CORAL/DAN 代表统计对齐，DANN/CDAN 代表对抗式对齐，DeepJDOT 代表最优传输思想，RCTA 则是本文的改进方法。

为降低无关变量干扰，所有方法共享同一 FCN 骨干和统一结果导出结构，但不同方法在训练轮数、对齐权重和模型选择指标上允许保留合理差异。当前 fixed-fold benchmark 中，基线方法训练轮数主要落在 84--120 之间，RCTA 训练轮数为 112；DeepJDOT 使用更大的 batch size 以支持 OT 代价计算，其余方法大体维持 batch size=32。这样的配置方式更符合 benchmark 的实际目的，即比较方法机制在合理调参后的表现，而非强行让所有方法使用完全一致但未必合适的超参数。

\subsection{运行命令与结果产物}
当前 benchmark 的实际运行入口为仓库中的批量脚本
\path{scripts/run_small_scale_round.sh}。对应的主实验命令为：

\begin{quote}
\small
\ttfamily
scripts/run\_small\_scale\_round.sh \textbackslash\\
--experiment-config \textbackslash\\
configs/experiment/\textbackslash\\
benchmark\_88\_2clusters\_11scenes\_8methods\_fixedfold.yaml
\end{quote}

运行结束后，脚手架会在 \path{runs/} 目录下统一生成每次运行的
\path{result.json}、\path{review.json}、checkpoint、分析特征文件以及图表；随后自动汇总到
\path{comparison_summary}
目录，形成整轮 benchmark 的 JSON/Markdown 对照表和汇总图件。也就是说，本章后续所有表格与图件理论上都可以从统一结果产物中回填，而不需要手工拼接零散结果。

\section{总体 benchmark 结果占位}
\subsection{8 个单源场景的结果表}
表~\ref{tab:benchmark-single-source} 预留了 8 个单源场景上的 benchmark 主结果位置。考虑到当前版本的论文重点是先完成结构化写作与实验回填接口，因此表中数值暂以占位形式呈现，待 benchmark 跑完后可直接回填目标域准确率或“准确率 / 平衡准确率”双指标。

\begin{table}[htbp]
\centering
\caption{8 个单源场景上的主结果（占位）}
\label{tab:benchmark-single-source}
\scriptsize
\resizebox{\textwidth}{!}{%
\begin{tabular}{lcccccccc}
\hline
场景 & Source Only & CORAL & DAN & DANN & CDAN & DeepJDOT & RCTA & Target Only \\
\hline
mode3$\rightarrow$mode6 & 待补充 & 待补充 & 待补充 & 待补充 & 待补充 & 待补充 & 待补充 & 待补充 \\
mode6$\rightarrow$mode3 & 待补充 & 待补充 & 待补充 & 待补充 & 待补充 & 待补充 & 待补充 & 待补充 \\
mode1$\rightarrow$mode2 & 待补充 & 待补充 & 待补充 & 待补充 & 待补充 & 待补充 & 待补充 & 待补充 \\
mode2$\rightarrow$mode1 & 待补充 & 待补充 & 待补充 & 待补充 & 待补充 & 待补充 & 待补充 & 待补充 \\
mode1$\rightarrow$mode5 & 待补充 & 待补充 & 待补充 & 待补充 & 待补充 & 待补充 & 待补充 & 待补充 \\
mode5$\rightarrow$mode1 & 待补充 & 待补充 & 待补充 & 待补充 & 待补充 & 待补充 & 待补充 & 待补充 \\
mode2$\rightarrow$mode5 & 待补充 & 待补充 & 待补充 & 待补充 & 待补充 & 待补充 & 待补充 & 待补充 \\
mode5$\rightarrow$mode2 & 待补充 & 待补充 & 待补充 & 待补充 & 待补充 & 待补充 & 待补充 & 待补充 \\
\hline
\end{tabular}
}
\normalsize
\end{table}

\subsection{3 个多源场景的结果表}
表~\ref{tab:benchmark-multi-source} 预留了 3 个多源场景上的主结果位置。该表的意义在于观察当前方法在“源域信息更丰富，但域间关系也更复杂”的设置下是否仍然能够保持稳定收益。

\begin{table}[htbp]
\centering
\caption{3 个多源场景上的主结果（占位）}
\label{tab:benchmark-multi-source}
\scriptsize
\resizebox{\textwidth}{!}{%
\begin{tabular}{lcccccccc}
\hline
场景 & Source Only & CORAL & DAN & DANN & CDAN & DeepJDOT & RCTA & Target Only \\
\hline
mode1+mode2$\rightarrow$mode5 & 待补充 & 待补充 & 待补充 & 待补充 & 待补充 & 待补充 & 待补充 & 待补充 \\
mode1+mode5$\rightarrow$mode2 & 待补充 & 待补充 & 待补充 & 待补充 & 待补充 & 待补充 & 待补充 & 待补充 \\
mode2+mode5$\rightarrow$mode1 & 待补充 & 待补充 & 待补充 & 待补充 & 待补充 & 待补充 & 待补充 & 待补充 \\
\hline
\end{tabular}
}
\normalsize
\end{table}

\subsection{方法总体汇总表}
为避免正文只罗列逐场景结果而缺少总体结论，表~\ref{tab:benchmark-summary} 预留了方法层面的均值汇总位置。建议在回填时至少包含：8 个单源场景平均准确率、3 个多源场景平均准确率、11 个场景总体平均准确率、总体平均平衡准确率以及相对 Source Only 的平均增益。

\begin{table}[htbp]
\centering
\caption{88-run benchmark 的方法总体汇总（占位）}
\label{tab:benchmark-summary}
\small
\resizebox{\textwidth}{!}{%
\begin{tabular}{lccccc}
\hline
方法 & 单源平均准确率 & 多源平均准确率 & 总体平均准确率 & 总体平均平衡准确率 & 相对 Source Only 平均增益 \\
\hline
Source Only & 待补充 & 待补充 & 待补充 & 待补充 & 0 \\
CORAL & 待补充 & 待补充 & 待补充 & 待补充 & 待补充 \\
DAN & 待补充 & 待补充 & 待补充 & 待补充 & 待补充 \\
DANN & 待补充 & 待补充 & 待补充 & 待补充 & 待补充 \\
CDAN & 待补充 & 待补充 & 待补充 & 待补充 & 待补充 \\
DeepJDOT & 待补充 & 待补充 & 待补充 & 待补充 & 待补充 \\
RCTA & 待补充 & 待补充 & 待补充 & 待补充 & 待补充 \\
Target Only & 待补充 & 待补充 & 待补充 & 待补充 & 待补充 \\
\hline
\end{tabular}
}
\end{table}

在最终版本中，本节将基于上述三张表从以下三个角度展开分析。第一，比较不同方法在总体平均性能上的高低，观察 RCTA 是否能够稳定优于 Source Only 和强基线。第二，区分单源场景和多源场景，判断方法收益是否具有场景依赖性。第三，利用 Target Only 这一监督参考观察当前无监督领域自适应方法距离“目标域可监督训练”还有多大差距，从而给出更有边界感的结论。

\section{定性分析图件占位}
除定量指标外，当前 benchmark 还会自动导出混淆矩阵、t-SNE 投影和汇总热力图。考虑到这些图件的最终版本依赖正式运行结果，本节先给出结构化占位，待 benchmark 跑完后可直接替换。

\begin{figure}[htbp]
\centering
\resultplaceholder[5.0cm]{11 个场景 $\times$ 8 种方法的 benchmark 热力图待补充}
\caption{88-run benchmark 的场景--方法热力图（占位）}
\label{fig:benchmark-heatmap}
\end{figure}

\begin{figure}[htbp]
\centering
\begin{minipage}{0.49\textwidth}
\centering
\resultplaceholder{代表性场景下强基线混淆矩阵待补充}
\par
\footnotesize (a) 强基线或 Source Only
\end{minipage}
\hfill
\begin{minipage}{0.49\textwidth}
\centering
\resultplaceholder{代表性场景下完整 RCTA 混淆矩阵待补充}
\par
\footnotesize (b) 完整 RCTA
\end{minipage}
\caption{代表性场景下基线方法与完整 RCTA 的混淆矩阵对比（占位）}
\label{fig:rcta-confusion}
\end{figure}

\begin{figure}[htbp]
\centering
\begin{minipage}{0.49\textwidth}
\centering
\resultplaceholder{代表性场景下强基线 t-SNE 图待补充}
\par
\footnotesize (a) 强基线或 Source Only
\end{minipage}
\hfill
\begin{minipage}{0.49\textwidth}
\centering
\resultplaceholder{代表性场景下完整 RCTA t-SNE 图待补充}
\par
\footnotesize (b) 完整 RCTA
\end{minipage}
\caption{代表性场景下基线方法与完整 RCTA 的域分布 t-SNE 对比（占位）}
\label{fig:rcta-tsne}
\end{figure}

在最终回填中，本节建议重点观察以下问题：误分类是否集中于少数易混类别；RCTA 是否能够提升混淆矩阵主对角线连续性；源域与目标域样本在 t-SNE 上是否形成更合理的跨域同类聚类；在多源场景中，多个源域与目标域之间是否出现更平滑的过渡结构。这些定性现象将用于支撑前文关于“稳定目标域利用”和“缓解类别混叠”的论证。

\section{模块递进验证的后续占位}
需要说明的是，当前论文周期内\textbf{必须完成的正式实验}是前述 88-run fixed-fold benchmark，而 A、A+B、A+B+C、A+B+C+D 形式的严格模块递进验证尚未被纳入当前主实验配置。因此，针对第三章至第六章各创新点的模块级消融，本节先保留结构化位置，待后续单独设计 ablation config 后再补入结果。

\begin{table}[htbp]
\centering
\caption{RCTA 模块递进验证结果（后续专用占位）}
\label{tab:rcta-ablation-plan}
\small
\begin{tabular}{lcccc}
\hline
模型组合 & mode2$\rightarrow$mode5 & mode3$\rightarrow$mode6 & mode5$\rightarrow$mode2 & mode6$\rightarrow$mode3 \\
\hline
A & 待补充 & 待补充 & 待补充 & 待补充 \\
A+B & 待补充 & 待补充 & 待补充 & 待补充 \\
A+B+C & 待补充 & 待补充 & 待补充 & 待补充 \\
A+B+C+D & 待补充 & 待补充 & 待补充 & 待补充 \\
\hline
\end{tabular}
\normalsize
\end{table}

保留这一占位的意义在于：正文的方法章节已经按创新点顺序展开，后续只要补上专门的 ablation 运行配置，就能够在不改论文结构的前提下形成从主 benchmark 到模块级验证的完整闭环。换言之，本文当前并非“没有消融位置”，而是“主 benchmark 先行、消融后补”。

\section{本章小结}
本章围绕当前实际需要执行的 88-run fixed-fold benchmark，系统说明了数据集规模、11 个场景的构成与折分配、统一实验协议、运行命令、主结果表和定性图件占位，并将严格模块消融作为后续专用占位单独说明。这样组织后，论文主实验与当前代码配置实现了一一对应，后续只需运行 benchmark 并把结果回填到本章既有结构中，即可形成完整实验论证。
